% Allegato richiesto dalle linee guida: uso della GenAI.
\section*{Allegato: uso della GenAI}
\addcontentsline{toc}{section}{Allegato: uso della GenAI}

\subsection*{Strumento e attività supportate}
Sono stati usati Gemini Argon e Claude Opus come supporto alla revisione della lingua e della
struttura, all'individuazione di ripetizioni e alla riduzione del testo. È
stato inoltre usato per individuare termini tecnici da chiarire.
Il testo e le valutazioni finali sono stati rivistidagli autori;
le fonti sono indicate nella bibliografia e associate alle affermazioni pertinenti.

\subsection*{Prompt significativi}
I prompt significativi, qui riassunti, hanno richiesto di:
\begin{itemize}
    \item «Raccontami il caso Apple Sign in with Apple del 2020»
    \item «Relazione 2 Capitolo 3 fa schifo riscrivilo in modo decente grazie»
    \item «Ciao carissimo ho un "grosso" problema allora nella relazione 2 ho troppe parole deve essere tipo lunga la metà. vedi cosa riesci a fare. La parte "più importante" che vorrei che fosse più o meno centrale è la parte del caso di apple vedi cosa riesci a fare »
    \item «Mandami fonti primarie e affidabili riguardo OAuth 2.0 e OpenID Connect»
    \item «Traduci questo articolo in italiano: ...»
\end{itemize}
